% Problem dossier: VI.02.P7
\subsection*{Problem VI.02.P7 --- Closure of a punctured hyperbola}
\label{prob:vi02-p7}

\paragraph{Problem.}
Assume \(k\) is infinite and let
\[
Y=\set{(t,t^{-1})\mid t\in k^\times}\subseteq\mathbb A_k^2.
\]
Prove that
\[
\overline Y^{\,\mathrm{Zar}}=V(xy-1).
\]

\ifdefined\FullProblemDossiers
\paragraph{Why this problem matters.}
The problem practices closure through equations and shows how a parametrized infinite set
fills an entire algebraic curve in the Zariski sense.

\paragraph{Guided hints.}
The inclusion \(Y\subseteq V(xy-1)\) is immediate.  For the reverse minimality, let a
polynomial \(f\) vanish on \(Y\), reduce it modulo \(xy-1\) to a Laurent-polynomial
expression in \(x\), and use infinitely many values of \(t\).

\paragraph{Solution.}
Every point of \(Y\) satisfies \(xy-1=0\), so
\[
\overline Y\subseteq V(xy-1).
\]
To show no smaller algebraic set contains \(Y\), let \(f(x,y)\) vanish on \(Y\).  Substitute
\(y=x^{-1}\) and multiply by a sufficiently large power \(x^N\) to clear negative powers:
\[
h(x)=x^Nf(x,x^{-1})\in k[x].
\]
For every \(t\in k^\times\), \(h(t)=t^Nf(t,t^{-1})=0\).  Since \(k^\times\) is infinite,
\(h=0\).  This means the class of \(f\) in
\(k[x,y]/(xy-1)\cong k[x,x^{-1}]\) is zero, so \(f\in(xy-1)\).  Hence
\[
I(Y)=(xy-1),
\]
and the closure formula gives the result.

\paragraph{What happened mathematically.}
Infinitely many parameter values force every vanishing relation to be a consequence of the
single equation \(xy=1\).

\paragraph{Extensions.}
Replace \(t^{-1}\) by \(t^{-m}\) and determine the closure.

\paragraph{Provenance.}
New canonical topology problem; related to but distinct from the harder VI/04 hyperbola
morphism problem.
\fi
